%% toc: P1 2 2.1
\section{Separation instances for relative complement and the intersection operation}
\label{sec:p01-w02b-p01-004-separation-instances}

Two constructions of \cref{sec:p01-w02-p01-005-set-operations} rest on an
instance of the separation schema: the relative complement \(A\setminus B\),
whose membership condition is clause 3 of
\cref{thm:p01-w02-p01-005-membership}, and the intersection operation on the
subsets of a set, which \cref{thm:p01-w02-p01-005-subset-algebra} asserts to
be a function from \(\mathcal P(S)\times\mathcal P(S)\) to \(\mathcal P(S)\).
This section writes both instances out in full: the \(\Lang_{\in}\)-formula, its
parameters, the bounding set, and the verification that the separated set has
the required members. For the intersection operation it also isolates the one
property of the bound that the construction needs, closure of the bound under
\(\cap\), and shows that the construction fails without it.

\emph{Inputs.} The language \(\Lang_{\in}\) and the conventions of
\cref{def:p01-w02-p01-005-membership}; extensionality, pairing and power set
from \cref{def:p01-w02-p01-005-existence-principles}; bounded set formation
(\cref{def:p01-w02-p01-005-bounded-formation}) and its existence and
uniqueness lemma \cref{thm:p01-w02-p01-005-separated-set}; the definitions of
\(A\cap B\) and \(A\setminus B\) in \cref{def:p01-w02-p01-005-operations}; the
membership condition of \(A\cap B\) (\cref{thm:p01-w02-p01-005-membership}(2));
the product bound \cref{thm:p01-w02-p01-005-product-bound}; the ordered pair
\(\langle a,b\rangle=\{\{a\},\{a,b\}\}\) and the ordered-pair identity of
\cref{sec:set-prerequisites-02}
\cite[ch.~``Sets, Relations, Functions, and Cardinality'', sec.~``Ordered pairs
and products'', definition ``Ordered pair'' and lemma ``Ordered-pair
identity'']{FrozenFormalLogicSetOperations}; and the definitions of binary
relation and function in \cref{sec:set-prerequisites-03}
\cite[ch.~``Sets, Relations, Functions, and Cardinality'', sec.~``Relations
and functions'', definition ``Function'']{FrozenFormalLogicSetOperations}.
\emph{Conclusions.} The membership condition and inclusion of the relative
complement obtained directly from separation
(\cref{thm:p01-w02b-p01-004-relative-complement}); \(\Lang_{\in}\)-formulas that
define the unordered pair, the ordered pair, binary intersection and the graph
of intersection (\cref{thm:p01-w02b-p01-004-defining-formulas}); for every set
\(Q\), the separated graph of intersection on \(Q\), its membership condition,
its single-valuedness, and the equivalence of its totality with closure of
\(Q\) under \(\cap\) (\cref{thm:p01-w02b-p01-004-intersection-graph}); and the
intersection operation on \(\mathcal P(S)\)
(\cref{thm:p01-w02b-p01-004-intersection-on-subsets}). \emph{Ambient
assumptions.} The logic is classical first-order logic with identity. No
principle beyond extensionality, empty set, pairing, union, power set and
separation is used; in particular neither replacement nor choice is used.

\subsection{Defining formulas}
\label{sec:p01-w02b-p01-004-formulas}

\begin{hbconvention}{def:p01-w02b-p01-004-variables}{Variables in the displayed formulas}
In each formula below the bound variables are distinct from one another and
from the free variables shown in its argument list. When variables are
substituted for the free variables of such a formula, the bound variables are
first renamed to variables that occur nowhere in the substituted variables, so
that no substituted variable is captured. The formulas are written with the
abbreviation \(x\notin a\) for \(\neg\,x\in a\) of
\cref{def:p01-w02-p01-005-membership}.
\end{hbconvention}

\begin{hbdefinition}{def:p01-w02b-p01-004-formulas}{Formulas for pairs, intersection and the intersection graph}
For variables \(w,z,a,b,u,X,Y,Q\) define the following formulas of
\(\Lang_{\in}\):
\begin{enumerate}
  \item the \emph{unordered-pair formula}
        \(\theta_{\mathrm{pair}}(w,a,b)\defeq
        \forall t\,\bigl(t\in w\leftrightarrow(t=a\lor t=b)\bigr)\);
  \item the \emph{ordered-pair formula}
        \(\theta_{\mathrm{op}}(z,a,b)\defeq
        \forall s\,\bigl(s\in z\leftrightarrow
        (\theta_{\mathrm{pair}}(s,a,a)\lor\theta_{\mathrm{pair}}(s,a,b))\bigr)\);
  \item the \emph{intersection formula}
        \(\theta_{\cap}(u,X,Y)\defeq
        \forall v\,\bigl(v\in u\leftrightarrow(v\in X\land v\in Y)\bigr)\);
  \item the \emph{intersection-graph formula}
        \[
          \begin{aligned}
            \gamma_{\cap}(w,Q)\defeq{}&
            \exists X\,\exists Y\,\exists p\,\exists u\,
            \bigl(X\in Q\land Y\in Q\land\theta_{\mathrm{op}}(p,X,Y)\\
            &\qquad\land\theta_{\cap}(u,X,Y)\land\theta_{\mathrm{op}}(w,p,u)\bigr).
          \end{aligned}
        \]
\end{enumerate}
Each is built from atomic formulas \(x\in y\) and \(x=y\) by connectives and
quantifiers only, so each is a formula of \(\Lang_{\in}\) with no defined
notation left to unfold. The free variables of \(\gamma_{\cap}(w,Q)\) are
exactly \(w\) and \(Q\).
\end{hbdefinition}

\begin{hblemma}{thm:p01-w02b-p01-004-defining-formulas}{proved=annex:p01-w02b-p01-004-defining-formulas}{The formulas define their intended sets}
For all sets \(w,z,a,b,u,X,Y,Q\):
\begin{enumerate}
  \item \(\theta_{\mathrm{pair}}(w,a,b)\leftrightarrow w=\{a,b\}\), and
        \(\theta_{\mathrm{pair}}(w,a,a)\leftrightarrow w=\{a\}\);
  \item \(\theta_{\mathrm{op}}(z,a,b)\leftrightarrow z=\langle a,b\rangle\);
  \item \(\theta_{\cap}(u,X,Y)\leftrightarrow u=X\cap Y\);
  \item \(\gamma_{\cap}(w,Q)\leftrightarrow\exists X\,\exists Y\,
        \bigl(X\in Q\land Y\in Q\land
        w=\langle\langle X,Y\rangle,X\cap Y\rangle\bigr)\).
\end{enumerate}
\end{hblemma}

\subsection{The relative complement}
\label{sec:p01-w02b-p01-004-relative-complement}

\begin{hblemma}{thm:p01-w02b-p01-004-relative-complement}{proved=annex:p01-w02b-p01-004-relative-complement}{The relative complement as a separated set}
Let \(A\) and \(B\) be sets. The formula \(x\notin B\) is a formula of
\(\Lang_{\in}\) with the single parameter \(B\), and the separation instance
for it, with bound \(A\), yields a unique set \(A\setminus B\) such that
\[
  \forall z\,\bigl(z\in A\setminus B\leftrightarrow(z\in A\land z\notin B)\bigr).
\]
Moreover \(A\setminus B\subseteq A\).
\end{hblemma}

The lemma is clause 3 of \cref{thm:p01-w02-p01-005-membership} together with
the inclusion \(A\setminus B\subseteq A\) of clause 6 of that lemma, each
obtained by one application of \cref{thm:p01-w02-p01-005-separated-set}.

\subsection{The graph of intersection}
\label{sec:p01-w02b-p01-004-graph}

\begin{hbdefinition}{def:p01-w02b-p01-004-graph}{The separated graph of intersection on a set}
Let \(Q\) be a set. By \cref{thm:p01-w02-p01-005-product-bound}, applied first
with \(A=B=Q\) and then with \(A=Q\times Q\) and \(B=Q\), the products
\(Q\times Q\) and \((Q\times Q)\times Q\) are sets. The \emph{separated graph
of intersection on \(Q\)} is
\[
  G^{Q}_{\cap}\defeq
  \{w\in(Q\times Q)\times Q:\gamma_{\cap}(w,Q)\}.
\]
The side conditions of \cref{def:p01-w02-p01-005-bounded-formation} hold: the
bound \((Q\times Q)\times Q\) is a set; the condition \(\gamma_{\cap}(w,Q)\) is
the \(\Lang_{\in}\)-formula of \cref{def:p01-w02b-p01-004-formulas} with the
single parameter \(Q\); and the variable for the formed set is chosen distinct
from \(w\), from \(Q\) and from every variable of \(\gamma_{\cap}\), so it is
not free in the condition.
\end{hbdefinition}

\begin{hbtheorem}{thm:p01-w02b-p01-004-intersection-graph}{proved=annex:p01-w02b-p01-004-intersection-graph}{Membership, single-valuedness and totality of the intersection graph}
Let \(Q\) be a set and write \(G=G^{Q}_{\cap}\).
\begin{enumerate}
  \item \(G\) is a set, \(G\subseteq(Q\times Q)\times Q\), so \(G\) is a binary
        relation from \(Q\times Q\) to \(Q\), and for every set \(w\)
        \[
          w\in G\leftrightarrow\exists X\,\exists Y\,
          \bigl(X\in Q\land Y\in Q\land X\cap Y\in Q\land
          w=\langle\langle X,Y\rangle,X\cap Y\rangle\bigr).
        \]
  \item If \(X\) and \(Y\) are sets and
        \(\langle\langle X,Y\rangle,u\rangle\in G\), then \(u=X\cap Y\).
        Consequently, if \(\langle m,u\rangle\in G\) and
        \(\langle m,u'\rangle\in G\), then \(u=u'\).
  \item \(G\) is a function from \(Q\times Q\) to \(Q\) iff
        \(X\cap Y\in Q\) for all \(X,Y\in Q\). In that case the value of \(G\)
        at \(\langle X,Y\rangle\) is \(X\cap Y\) for all \(X,Y\in Q\).
\end{enumerate}
\end{hbtheorem}

\begin{hbcorollary}{thm:p01-w02b-p01-004-intersection-on-subsets}{proved=annex:p01-w02b-p01-004-intersection-on-subsets}{Intersection is an operation on the subsets of a set}
Let \(S\) be a set and \(Q=\mathcal P(S)\). Then \(X\cap Y\in Q\) for all
\(X,Y\in Q\), and \(G^{Q}_{\cap}\) is a function from \(Q\times Q\) to \(Q\)
whose value at \(\langle X,Y\rangle\) is \(X\cap Y\). Its members are exactly
the triples \(\langle\langle X,Y\rangle,X\cap Y\rangle\) with
\(X,Y\subseteq S\). This function is the intersection operation \(G_{\cap}\)
of \cref{thm:p01-w02-p01-005-subset-algebra}.
\end{hbcorollary}

\subsection{Worked example}
\label{sec:p01-w02b-p01-004-example}

\begin{hbexample}{def:p01-w02b-p01-004-example}{The intersection graph on the subsets of a singleton}
Let \(e_0=\varnothing\) and \(e_1=\{\varnothing\}\), as in
\cref{def:p01-w02-p01-005-example}, and let \(S=e_1\), \(Q=\mathcal P(S)\).

\emph{The carrier.} A set \(X\) lies in \(Q\) iff every member of \(X\) is a
member of \(e_1\), that is, equals \(\varnothing\). If such an \(X\) has a
member, that member is \(\varnothing\), so \(\varnothing\in X\); then \(X\) and
\(e_1\) have the same members and \(X=e_1\) by extensionality. If \(X\) has no
member, then \(X=e_0\) by extensionality. Both \(e_0\subseteq S\) and
\(e_1\subseteq S\) hold. So \(z\in Q\leftrightarrow(z=e_0\lor z=e_1)\), and
\(Q=\{e_0,e_1\}\).

\emph{The intersections.} By \cref{thm:p01-w02-p01-005-membership}(2), no set
is a member of \(e_0\cap Y\) or of \(X\cap e_0\), so both equal \(e_0\); and
\(e_1\cap e_1=e_1\) by idempotence (\cref{thm:p01-w02-p01-005-laws}(1)).

\emph{The graph.} By \cref{thm:p01-w02b-p01-004-intersection-on-subsets}, the
members of \(G^{Q}_{\cap}\) are exactly
\[
  \langle\langle e_0,e_0\rangle,e_0\rangle,\quad
  \langle\langle e_0,e_1\rangle,e_0\rangle,\quad
  \langle\langle e_1,e_0\rangle,e_0\rangle,\quad
  \langle\langle e_1,e_1\rangle,e_1\rangle .
\]
Each is a member of \((Q\times Q)\times Q\) because its last coordinate lies in
\(Q\), and each satisfies \(\gamma_{\cap}(w,Q)\) with the witnesses \(X,Y\)
read off from its first coordinate.

\emph{The complements.} By \cref{thm:p01-w02b-p01-004-relative-complement},
\(z\in S\setminus e_0\) iff \(z\in e_1\land z\notin\varnothing\), that is, iff
\(z\in e_1\); so \(S\setminus e_0=e_1\). And \(z\in S\setminus e_1\) iff
\(z\in e_1\land z\notin e_1\), which no \(z\) satisfies; so
\(S\setminus e_1=e_0\). Complement exchanges the bottom \(e_0\) and the top
\(e_1\) of the two-element Boolean algebra \(\mathcal P(e_1)\).
\end{hbexample}

\subsection{Failure when the bound is not closed}
\label{sec:p01-w02b-p01-004-failure}

\begin{hbwarning}{def:p01-w02b-p01-004-closure-warning}{The separation formula alone does not make an operation}
The formula \(\gamma_{\cap}(w,Q)\) describes the intended triples for every
set \(Q\), but the separated set \(G^{Q}_{\cap}\) keeps only those triples
whose last coordinate lies in the bound. Let \(e_1=\{\varnothing\}\) and
\(e_2=\{\{\varnothing\}\}\), as in \cref{def:p01-w02-p01-005-example}, and let
\(Q=\{e_1,e_2\}\).
\begin{itemize}
  \item \(e_1\cap e_2=\varnothing\). The only member of \(e_1\) is
        \(\varnothing\), and \(\varnothing\notin e_2\), because the only member
        of \(e_2\) is \(\{\varnothing\}\), which differs from \(\varnothing\);
        so by \cref{thm:p01-w02-p01-005-membership}(2) the intersection has no
        member.
  \item \(\varnothing\notin Q\). The set \(\varnothing\) differs from \(e_1\),
        which has the member \(\varnothing\), and from \(e_2\), which has the
        member \(e_1\).
\end{itemize}
So \(Q\) is not closed under \(\cap\), and by
\cref{thm:p01-w02b-p01-004-intersection-graph}(3) the relation
\(G^{Q}_{\cap}\) is not a function from \(Q\times Q\) to \(Q\): by clause 2 of
that theorem the only value it could take at \(\langle e_1,e_2\rangle\) is
\(\varnothing\), and by clause 1 no triple with last coordinate \(\varnothing\)
is a member. A proof that a separated graph is an operation on \(Q\) must
therefore show that every value lies in \(Q\); for \(Q=\mathcal P(S)\) this is
the step \(X\cap Y\subseteq S\) of
\cref{thm:p01-w02b-p01-004-intersection-on-subsets}.
\end{hbwarning}

\subsection{Relationships}
\label{sec:p01-w02b-p01-004-relationships}

\Cref{thm:p01-w02b-p01-004-relative-complement} supplies the separation step
behind clause 3 of \cref{thm:p01-w02-p01-005-membership}, and through it the
laws of \cref{thm:p01-w02-p01-005-laws} that mention a relative complement.
\Cref{thm:p01-w02b-p01-004-intersection-on-subsets} supplies the intersection
operation \(G_{\cap}\) of \cref{thm:p01-w02-p01-005-subset-algebra}, and hence
the meet of the Boolean algebra of subsets used with filters and ultrafilters in
\cref{def:p01-bpi-filter-ideal}. The formulas \(\theta_{\mathrm{pair}}\) and
\(\theta_{\mathrm{op}}\) are the unfolding of unordered and ordered pairs used
in the proof of \cref{thm:p01-w02-p01-005-product-bound}, written as named
formulas. \Cref{thm:p01-w02b-p01-004-intersection-graph} is an instance of the
procedure \cref{def:p01-w02-p01-005-bounding-procedure}: its clause 3 states
exactly when the first step of that procedure, finding a bound that contains
every intended member, succeeds for the bound \((Q\times Q)\times Q\).

\subsection{Sources}
\label{sec:p01-w02b-p01-004-sources}

The language \(\Lang_{\in}\), the axioms of extensionality, pairing and power
set, the separation schema, the ordered pair and the ordered-pair identity, and
the definitions of binary relation and of function are taken from the
frozen manuscript \cite{FrozenFormalLogicSetOperations}, at the locators given above
and in \cref{sec:p01-w02-p01-005-set-operations}. The defining formulas, the
separated graph of intersection on a set, the lemmas, the theorem, the
corollary, the example, the counterexample and all proofs are editorial
completions from those principles and are not attributed to an external
presentation.
